\begin{frame}{Phương pháp Nhân phức (Complex Multiplication)}
  \begin{block}{Ý tưởng cốt lõi}
    Thay vì chọn $(a,b)$ ngẫu nhiên rồi đếm điểm,\\
    CM \emph{đảo chiều}: bắt đầu từ bậc nhóm $r$ mong muốn $\Rightarrow$ suy ra đường cong
  \end{block}
  \vspace{0.3em}
  \textbf{Điều kiện CM:} Tồn tại đường cong trên $\mathbb{F}_p$ với vết Frobenius $t$ khi:
  \[
    4p = t^2 + |D| \cdot y^2 \quad (D < 0 \text{ là biệt thức CM})
  \]
  \vspace{0.2em}
  \textbf{Quy trình 4 bước:}
  \begin{enumerate}
    \item Tính $j$-invariant từ đa thức lớp Hilbert $H_D(x) \pmod{p}$
    \item Suy ra hệ số đường cong $(a, b)$ từ $j$
    \item Kiểm tra xoắn (twist test): chọn đường cong có $r \mid \#E$
    \item Tìm điểm sinh $G$ bậc $r$ bằng nhân cofactor
  \end{enumerate}
  \vspace{0.2em}
  Với $D = -3$: $H_{-3}(x) = x \Rightarrow j = 0 \Rightarrow y^2 = x^3 + b$ (dạng đơn giản nhất)
\end{frame}
